% Generated draft: requires mathematical and source review.
\documentclass[11pt]{article}
\usepackage[margin=25mm]{geometry}
\usepackage{fontspec}
\setmainfont{Latin Modern Roman}
\usepackage{amsmath,amssymb,mathtools,unicode-math}
\setmathfont{Latin Modern Math}
\usepackage{enumitem,needspace,xurl,hyperref}
\setlength{\parindent}{0pt}
\setlength{\parskip}{5pt}
\setlength{\emergencystretch}{3em}
\newcommand{\N}{\mathbb{N}}
\newcommand{\Z}{\mathbb{Z}}
\newcommand{\Q}{\mathbb{Q}}
\newcommand{\R}{\mathbb{R}}
\newcommand{\C}{\mathbb{C}}
\newcommand{\F}{\mathbb{F}}
\newcommand{\A}{\mathbb{A}}
\newcommand{\PP}{\mathbb{P}}
\newcommand{\E}{\mathbb{E}}
\newcommand{\eps}{\varepsilon}
\newcommand{\dd}{\,\mathrm{d}}
\begin{document}
\section*{UnsolvedMath ID: 701667}
\textbf{Source problem code:} catalogue entry
\subsection{Definitions and mathematical statement}
Let \(k\) be a field and \(F\) a free group, with no restriction imposed on the characteristic of \(k\) or the rank of \(F\). The group algebra \(k[F]\) consists of finite sums
\[
u=\sum_{g\in F}c_g g,\qquad c_g\in k,
\]
with multiplication induced by the group operation and extended \(k\)-bilinearly. Define
\[
\operatorname{supp}(u)=\{g\in F:c_g\ne0\}.
\]
The two-sided ideal generated by \(u\) is
\[
I(u)=\left\{\sum_{i=1}^{m}a_iub_i:m\in\Z_{\ge0},\ a_i,b_i\in k[F]\right\},
\]
where the empty sum is zero.

The question proposed by G. M. Bergman is whether, for every field \(k\), every free group \(F\), and every \(u\in k[F]\),
\[
|\operatorname{supp}(u)|>1
\quad\Longrightarrow\quad
I(u)\ne k[F].
\]
Equivalently, under these hypotheses, must the identity element \(1_{k[F]}\) lie outside \(I(u)\)? The supplied registry labels this question open; that status is not independently verified here.
\subsection{Short English statement}
Must every element supported on at least two distinct group elements in the group algebra of any free group over any field generate a proper two-sided ideal? This is the question attributed to G. M. Bergman in Kourovka Notebook Problem 20.9.
\subsection{Sources}
\begin{itemize}
\item Local UnsolvedMath registry, permanent ID \(701667\), source problem code \texttt{catalogue entry}, title \emph{Kourovka Notebook Problem 20.9}. The supplied record attributes the problem to G. M. Bergman, gives the proposed year as 2022, and labels its status open.
\item \emph{The Kourovka Notebook (Unsolved Problems in Group Theory)}, Problem 20.9; the supplied background cites the 20th issue (2022) and the 21st issue edition (2026), arXiv:1401.0300v48. The supplied collection description names E. I. Khukhro and V. D. Mazurov as editors.
\item G. M. Bergman, \emph{Commun. Algebra} \textbf{49}, no. 9 (2021), 3760--3776. The supplied statement identifies this question as one of several related questions in that article; it does not supply the article title.
\end{itemize}
\end{document}
